% ============================================================
% PHẦN 3.7 — CHIẾN LƯỢC THU HÚT KHÁCH HÀNG ĐẦU TIÊN
% ------------------------------------------------------------
% Thuộc: PA3-02
% Người phụ trách chính: Nguyễn Anh Thư – CEO/Project Lead
% Phối hợp: Trần Minh Quang
%
% CÂU HỎI CẦN TRẢ LỜI:
%   - Khách hàng đầu tiên sẽ được tìm từ đâu?
%   - Nhóm sẽ cung cấp ưu đãi gì để giảm rủi ro cho khách hàng?
%   - Điều kiện để một đơn vị được tham gia pilot là gì?
%   - Quy trình từ tiếp cận đến pilot gồm những bước nào?
%   - Sau pilot, điều kiện nào khiến khách hàng chuyển sang trả phí?
%
% OUTPUT BẮT BUỘC:
%   Chiến lược tìm khách hàng đầu tiên, gồm:
%     - Tiêu chí chọn khách hàng
%     - Nguồn danh sách khách hàng
%     - Cách liên hệ
%     - Đề nghị pilot
%     - Hình thức ưu đãi
%     - Điều kiện chuyển sang trả phí
%
% TIÊU CHÍ HOÀN THÀNH:
%   - Chiến lược phải chỉ ra cách tìm khách hàng THẬT, không chỉ ghi "chạy quảng cáo"
%   - Không khẳng định hiệu quả doanh thu khi chưa có pilot
% ============================================================

% TODO (Nguyễn Anh Thư): Viết nội dung mục 3.7
% Đây là phần quan trọng nhất của PA3 — phải cụ thể và khả thi với nguồn lực nhóm

\subsection{Chiến lược thu hút khách hàng đầu tiên}

Trong giai đoạn MVP, 1Stream không ưu tiên tìm khách hàng đầu tiên bằng quảng cáo đại trà. Sản phẩm vẫn cần nhóm hỗ trợ trực tiếp trong quá trình chuẩn hóa dữ liệu, cấu hình nội dung, kiểm thử phản hồi và theo dõi phiên livestream. Vì vậy, chiến lược phù hợp hơn là tiếp cận có chọn lọc từng đơn vị, ưu tiên các mối quan hệ có mức độ tin cậy ban đầu và chỉ mời pilot những khách hàng đáp ứng đủ điều kiện vận hành.

Mục tiêu của chiến lược này không phải thu được số lượng lớn lượt đăng ký, mà là tìm được một nhóm nhỏ khách hàng thật có vấn đề rõ ràng, sẵn sàng cung cấp dữ liệu, tham gia thử nghiệm và đưa ra quyết định tiếp tục hoặc dừng sau khi pilot kết thúc.

Trong 90 ngày đầu, nhóm đặt mục tiêu:

\begin{itemize}
    \item xây dựng danh sách tối thiểu 60 đơn vị tiềm năng;
    \item sàng lọc được khoảng 30 đơn vị phù hợp với ICP;
    \item tiếp cận cá nhân hóa tối thiểu 24 đơn vị;
    \item hoàn thành ít nhất 5 cuộc trao đổi hoặc demo;
    \item triển khai pilot với khoảng 3 đơn vị;
    \item chuyển đổi tối thiểu 1 đơn vị sang hình thức trả phí hoặc đặt cọc cho giai đoạn tiếp theo.
\end{itemize}

Các con số trên là mục tiêu thử nghiệm phục vụ quá trình học từ thị trường, không phải kết quả đã đạt được hoặc cam kết thương mại của dự án.

\subsubsection{Nguyên tắc lựa chọn khách hàng đầu tiên}

Khách hàng đầu tiên cần đồng thời đáp ứng hai yêu cầu: có vấn đề đủ rõ để cần giải pháp và có điều kiện phối hợp để nhóm triển khai pilot trong phạm vi nguồn lực hiện tại.

Nhóm không lựa chọn khách hàng chỉ dựa trên quy mô thương hiệu hoặc số lượng người theo dõi. Một trung tâm nhỏ nhưng có quy trình tuyển sinh rõ ràng, dữ liệu sẵn có và người phụ trách tích cực có thể phù hợp hơn một chuỗi trung tâm lớn với quy trình phê duyệt phức tạp.

\begin{table}[htbp]
    \centering
    \caption{Tiêu chí lựa chọn khách hàng đầu tiên của 1Stream}
    \label{tab:first-customer-selection-criteria}
    \begin{tabular}{
        p{0.25\textwidth}
        p{0.50\textwidth}
        p{0.15\textwidth}
    }
        \hline
        \textbf{Tiêu chí}
        & \textbf{Dấu hiệu phù hợp}
        & \textbf{Mức ưu tiên} \\
        \hline

        Thuộc phân khúc mục tiêu
        &
        Là trung tâm ngoại ngữ hoặc đơn vị bán khóa học có hoạt động tuyển sinh trên Facebook hoặc TikTok.
        &
        Bắt buộc \\

        \hline

        Có nhu cầu livestream
        &
        Đang livestream định kỳ, từng tổ chức livestream hoặc có kế hoạch triển khai chiến dịch livestream trong thời gian gần.
        &
        Cao \\

        \hline

        Đội ngũ vận hành còn mỏng
        &
        Đội marketing, nội dung và tư vấn có quy mô nhỏ; một người phải đảm nhận nhiều công việc.
        &
        Cao \\

        \hline

        Có dữ liệu tương đối chuẩn hóa
        &
        Có học phí, lịch học, chương trình học, ưu đãi, câu hỏi thường gặp, hình ảnh hoặc video được phép sử dụng.
        &
        Bắt buộc \\

        \hline

        Có quy trình thu lead
        &
        Đã sử dụng inbox, biểu mẫu, bảng tính hoặc nhân viên tư vấn để tiếp nhận khách hàng sau livestream.
        &
        Cao \\

        \hline

        Tiếp cận được người quyết định
        &
        Có thể trao đổi với chủ trung tâm, giám đốc, trưởng marketing hoặc người có quyền phê duyệt pilot.
        &
        Cao \\

        \hline

        Sẵn sàng phối hợp
        &
        Có người phụ trách cung cấp dữ liệu, duyệt nội dung, tham gia phản hồi và xử lý các tình huống phát sinh.
        &
        Bắt buộc \\

        \hline

        Có khả năng chi trả
        &
        Đang chi ngân sách cho host, nội dung, quảng cáo, nhân viên trực bình luận hoặc các công cụ hỗ trợ tương tự.
        &
        Trung bình đến cao \\

        \hline
    \end{tabular}
\end{table}

Các đơn vị sau chưa được ưu tiên trong giai đoạn đầu:

\begin{itemize}
    \item chưa từng livestream và chưa có kế hoạch triển khai cụ thể;
    \item không có dữ liệu khóa học hoặc dữ liệu thường xuyên thay đổi nhưng không có người kiểm soát;
    \item yêu cầu tích hợp CRM, ERP, SLA hoặc tùy biến sâu ngay từ đầu;
    \item kỳ vọng 1Stream thay thế hoàn toàn host và nhân viên tuyển sinh;
    \item yêu cầu cam kết về doanh thu hoặc số lượng lead khi chưa chạy thử nghiệm;
    \item không có quyền sử dụng hình ảnh, giọng nói, video hoặc tài khoản nền tảng;
    \item không bố trí được người duyệt thông tin và phối hợp trong quá trình pilot.
\end{itemize}

\subsubsection{Nguồn tìm kiếm danh sách khách hàng}

Danh sách khách hàng đầu tiên được xây dựng từ các nguồn thực tế mà nhóm có thể chủ động tiếp cận, thay vì phụ thuộc vào quảng cáo trả phí.

\begin{table}[htbp]
    \centering
    \caption{Nguồn tìm kiếm khách hàng đầu tiên của 1Stream}
    \label{tab:first-customer-sources}
    \begin{tabular}{
        p{0.21\textwidth}
        p{0.39\textwidth}
        p{0.29\textwidth}
    }
        \hline
        \textbf{Nguồn}
        & \textbf{Cách thực hiện}
        & \textbf{Thông tin cần ghi nhận} \\
        \hline

        Mạng lưới cá nhân của nhóm
        &
        Nhờ giảng viên, mentor, bạn bè, cựu sinh viên, giáo viên và người quen giới thiệu chủ hoặc nhân viên trung tâm ngoại ngữ.
        &
        Người giới thiệu, vai trò người liên hệ, mức độ quan hệ và nhu cầu đã biết. \\

        \hline

        Google Maps và công cụ tìm kiếm
        &
        Tìm các trung tâm ngoại ngữ tại TP.HCM theo khu vực, kiểm tra website, fanpage và thông tin liên hệ.
        &
        Quy mô, địa điểm, kênh hoạt động, số điện thoại, email và tên người phụ trách nếu có. \\

        \hline

        Facebook và TikTok
        &
        Tìm các fanpage hoặc tài khoản thường đăng nội dung tuyển sinh, video ngắn, livestream chữa đề hoặc tư vấn khóa học.
        &
        Tần suất nội dung, phiên livestream gần nhất, lượng tương tác và các câu hỏi phổ biến. \\

        \hline

        Cộng đồng ngành giáo dục
        &
        Quan sát các nhóm dành cho chủ trung tâm, giáo viên, marketing giáo dục, tuyển sinh và vận hành trung tâm.
        &
        Vấn đề được thảo luận, người có vai trò quản lý và cơ hội xin giới thiệu. \\

        \hline

        Đối tác cung cấp dịch vụ
        &
        Tiếp cận agency quảng cáo giáo dục, đơn vị quay dựng, chatbot, CRM, website hoặc tư vấn vận hành trung tâm.
        &
        Khả năng giới thiệu khách hàng, phạm vi hợp tác và cơ chế ghi nhận nguồn lead. \\

        \hline

        Sự kiện và workshop
        &
        Tham gia các buổi chia sẻ về giáo dục, tuyển sinh, marketing hoặc ứng dụng AI để kết nối trực tiếp.
        &
        Danh thiếp, vai trò, nhu cầu, thời điểm phù hợp để liên hệ lại. \\

        \hline
    \end{tabular}
\end{table}

Mỗi đơn vị được lưu trong một bảng theo dõi chung với các trường tối thiểu sau:

\begin{itemize}
    \item tên trung tâm;
    \item website, fanpage hoặc tài khoản TikTok;
    \item địa chỉ và phạm vi hoạt động;
    \item người liên hệ và vai trò;
    \item nguồn tìm thấy hoặc người giới thiệu;
    \item tần suất livestream ước tính;
    \item vấn đề quan sát được;
    \item mức độ đầy đủ của dữ liệu;
    \item điểm phù hợp ICP;
    \item trạng thái tiếp cận;
    \item bước tiếp theo;
    \item người phụ trách trong nhóm.
\end{itemize}

\subsubsection{Thứ tự ưu tiên tiếp cận}

Để tăng xác suất nhận được phản hồi và giảm chi phí xây dựng niềm tin, nhóm tiếp cận khách hàng theo thứ tự sau:

\begin{enumerate}
    \item khách hàng được người quen, giảng viên, mentor hoặc đối tác giới thiệu;
    \item trung tâm đã từng livestream và nhóm xác định được vấn đề cụ thể;
    \item trung tâm đang tuyển nhân sự livestream, marketing hoặc tư vấn tuyển sinh;
    \item trung tâm hoạt động tích cực trên Facebook hoặc TikTok nhưng chưa duy trì livestream ổn định;
    \item danh sách tiếp cận lạnh từ Google Maps, website và các cộng đồng ngành.
\end{enumerate}

Các mối quan hệ được giới thiệu được ưu tiên vì người giới thiệu giúp giảm rủi ro cảm nhận và tạo điều kiện để nhóm tiếp cận đúng người có quyền quyết định. Tuy nhiên, khách hàng được giới thiệu vẫn phải đáp ứng các điều kiện pilot, không được tự động lựa chọn chỉ vì có quan hệ quen biết.

\subsubsection{Cách liên hệ khách hàng}

Nhóm sử dụng cách tiếp cận cá nhân hóa theo từng tài khoản. Trước khi liên hệ, thành viên phụ trách cần xem xét hoạt động gần nhất của trung tâm và ghi nhận ít nhất một vấn đề hoặc cơ hội cụ thể.

Nội dung liên hệ ban đầu không bắt đầu bằng danh sách tính năng. Tin nhắn cần thể hiện rằng nhóm đã tìm hiểu hoạt động của trung tâm, nêu một vấn đề có thể quan sát và đề nghị một bước tiếp theo có mức cam kết thấp.

Ví dụ về tin nhắn tiếp cận:

\begin{quote}
Chào anh/chị, nhóm em thấy trung tâm đang có các nội dung tuyển sinh cho khóa học [tên khóa học] trên [Facebook/TikTok] và người xem thường hỏi về [học phí, lịch học hoặc đầu vào].

Nhóm em đang phát triển 1Stream, một giải pháp hỗ trợ trung tâm chuẩn bị nội dung livestream, phản hồi các câu hỏi phổ biến từ dữ liệu đã được duyệt và chuyển những trường hợp cần tư vấn sâu cho nhân viên.

Nhóm em có thể hỗ trợ kiểm tra miễn phí một phiên livestream hoặc quy trình hiện tại và gửi lại bản tóm tắt các điểm có nguy cơ làm tăng công việc vận hành hoặc bỏ sót khách hàng tiềm năng. Anh/chị có thể dành khoảng 20 phút để trao đổi trong tuần này không?
\end{quote}

Các kênh liên hệ được lựa chọn tùy theo thông tin có sẵn:

\begin{itemize}
    \item nhắn tin trực tiếp qua fanpage hoặc tài khoản cá nhân của người phụ trách;
    \item email đến chủ trung tâm, trưởng marketing hoặc bộ phận tuyển sinh;
    \item Zalo hoặc điện thoại nếu thông tin được công khai cho mục đích kinh doanh;
    \item liên hệ thông qua người giới thiệu;
    \item gặp trực tiếp tại sự kiện hoặc sau khi đã có lịch hẹn.
\end{itemize}

Nhóm không gửi cùng một nội dung đại trà đến số lượng lớn trung tâm. Mỗi tin nhắn cần được điều chỉnh theo loại khóa học, kênh đang sử dụng và vấn đề quan sát được của đơn vị.

\subsubsection{Chuỗi tiếp cận khách hàng}

Một khách hàng có thể được liên hệ tối đa bốn lần trong khoảng 14 ngày.

\begin{table}[htbp]
    \centering
    \caption{Chuỗi tiếp cận khách hàng đầu tiên}
    \label{tab:first-customer-outreach-sequence}
    \begin{tabular}{
        p{0.14\textwidth}
        p{0.26\textwidth}
        p{0.48\textwidth}
    }
        \hline
        \textbf{Thời điểm}
        & \textbf{Mục tiêu}
        & \textbf{Nội dung} \\
        \hline

        Ngày 1
        &
        Tạo cuộc trao đổi ban đầu
        &
        Gửi tin nhắn cá nhân hóa, nêu vấn đề quan sát được và đề nghị audit hoặc trao đổi ngắn. \\

        \hline

        Ngày 3--4
        &
        Cung cấp thêm giá trị
        &
        Gửi checklist, ví dụ về nhóm câu hỏi lặp lại hoặc một nhận xét ngắn liên quan đến phiên livestream gần nhất. \\

        \hline

        Ngày 7
        &
        Mời thực hiện hành động
        &
        Mời đặt lịch demo 20--30 phút hoặc nhận bản audit ngắn. \\

        \hline

        Ngày 12--14
        &
        Kết thúc vòng liên hệ
        &
        Gửi tin nhắn cuối, xác nhận sẽ không tiếp tục làm phiền nếu thời điểm chưa phù hợp và xin phép liên hệ lại trong tương lai. \\

        \hline
    \end{tabular}
\end{table}

Nếu khách hàng không phản hồi sau chuỗi trên, đơn vị được chuyển sang trạng thái theo dõi. Nhóm không tiếp tục nhắn tin thường xuyên, trừ khi xuất hiện tín hiệu mới như trung tâm bắt đầu livestream, tuyển nhân sự hoặc triển khai một chiến dịch tuyển sinh mới.

\subsubsection{Chương trình khách hàng đồng hành ban đầu}

Nhóm xây dựng chương trình \textbf{1Stream Early Partner Pilot} dành cho một số trung tâm đáp ứng tiêu chí ICP và sẵn sàng phối hợp kiểm chứng sản phẩm.

Chương trình không được quảng bá như một phiên bản miễn phí đầy đủ của sản phẩm. Đây là hoạt động thử nghiệm có giới hạn, trong đó hai bên cùng thống nhất mục tiêu, dữ liệu, trách nhiệm và tiêu chí đánh giá.

\paragraph{Quyền lợi của khách hàng tham gia}

\begin{itemize}
    \item được audit miễn phí một phiên livestream hoặc quy trình vận hành hiện tại;
    \item được cung cấp mẫu chuẩn hóa dữ liệu khóa học và câu hỏi thường gặp;
    \item được nhóm hỗ trợ trực tiếp trong quá trình onboarding;
    \item được tạo một số nội dung hoặc cấu hình thử nghiệm trong phạm vi đã thống nhất;
    \item được hỗ trợ kiểm thử trước khi triển khai phiên pilot;
    \item được nhận báo cáo kết quả sau pilot;
    \item được ưu tiên phản hồi và sửa các lỗi quan trọng thuộc phạm vi MVP;
    \item được hưởng mức phí pilot ưu đãi hoặc được khấu trừ khoản đặt cọc khi chuyển sang gói sử dụng tiếp theo.
\end{itemize}

\paragraph{Cam kết của khách hàng}

\begin{itemize}
    \item cung cấp dữ liệu đúng, cập nhật và có quyền sử dụng;
    \item chỉ định ít nhất một người phụ trách phối hợp;
    \item duyệt nội dung, dữ liệu và quy tắc phản hồi trước khi vận hành;
    \item tham gia buổi onboarding và buổi đánh giá sau pilot;
    \item thông báo kịp thời khi học phí, lịch học hoặc chính sách thay đổi;
    \item không yêu cầu các tính năng ngoài phạm vi MVP đã thống nhất;
    \item cho phép nhóm sử dụng số liệu tổng hợp và đã ẩn danh để đánh giá sản phẩm;
    \item cân nhắc cho phép sử dụng phản hồi hoặc case study nếu pilot tạo ra giá trị.
\end{itemize}

\subsubsection{Phạm vi đề nghị pilot}

Pilot được thu hẹp để giảm rủi ro cho khách hàng và bảo đảm nhóm có thể hỗ trợ tốt.

\begin{table}[htbp]
    \centering
    \caption{Phạm vi đề nghị pilot của 1Stream}
    \label{tab:first-customer-pilot-offer}
    \begin{tabular}{
        p{0.23\textwidth}
        p{0.56\textwidth}
    }
        \hline
        \textbf{Hạng mục}
        & \textbf{Phạm vi đề xuất} \\
        \hline

        Sản phẩm hoặc khóa học
        &
        Một khóa học hoặc một chương trình tuyển sinh cụ thể. \\

        \hline

        Nền tảng
        &
        Một nền tảng mà MVP đang hỗ trợ ổn định, ưu tiên Facebook hoặc TikTok tùy khả năng kỹ thuật thực tế. \\

        \hline

        Dữ liệu
        &
        Học phí, lịch học, thông tin đầu vào, ưu đãi, câu hỏi thường gặp và nội dung đã được phê duyệt. \\

        \hline

        Nội dung
        &
        Một số kịch bản, video hoặc nội dung mẫu trong giới hạn đã thống nhất. \\

        \hline

        Phản hồi
        &
        Chỉ xử lý các câu hỏi thuộc nhóm dữ liệu đã được cấu hình; trường hợp ngoài phạm vi được chuyển cho nhân viên. \\

        \hline

        Thời gian
        &
        Từ 14 đến 30 ngày, tùy mức độ sẵn sàng của dữ liệu và nền tảng. \\

        \hline

        Số phiên
        &
        Một hoặc một số phiên thử nghiệm giới hạn, không cam kết vận hành liên tục 24/7. \\

        \hline

        Hỗ trợ
        &
        Hỗ trợ onboarding, kiểm thử, theo dõi phiên và đánh giá kết quả. \\

        \hline

        Kết quả
        &
        Báo cáo về thời gian triển khai, câu hỏi được xử lý, trường hợp chuyển người thật, lead được ghi nhận, lỗi phát sinh và phản hồi của người dùng. \\

        \hline
    \end{tabular}
\end{table}

\subsubsection{Hình thức ưu đãi và giảm rủi ro}

Ưu đãi được thiết kế để giảm rủi ro khi khách hàng thử một sản phẩm mới, nhưng vẫn cần kiểm chứng mức sẵn sàng trả phí.

\begin{enumerate}
    \item \textbf{Audit miễn phí:}

    Nhóm phân tích một phiên livestream, một tập bình luận hoặc quy trình vận hành hiện tại và đưa ra nhận xét ban đầu. Audit giúp khách hàng nhận được giá trị trước khi quyết định tham gia pilot.

    \item \textbf{Onboarding có hỗ trợ:}

    Nhóm trực tiếp hỗ trợ chuẩn hóa dữ liệu, thiết lập câu hỏi thường gặp và kiểm thử các trường hợp mẫu. Khách hàng không phải tự cấu hình toàn bộ hệ thống.

    \item \textbf{Phạm vi thử nghiệm nhỏ:}

    Pilot chỉ triển khai trên một khóa học, một kênh và một giai đoạn ngắn. Cách này giúp khách hàng hạn chế ảnh hưởng đến hoạt động chính.

    \item \textbf{Quy trình duyệt trước:}

    Khách hàng được kiểm tra dữ liệu, nội dung và phản hồi mẫu trước khi hệ thống được sử dụng trong phiên thử nghiệm.

    \item \textbf{Quyền dừng thử nghiệm:}

    Khách hàng có thể yêu cầu tạm dừng khi phát hiện thông tin sai, nội dung không phù hợp hoặc rủi ro đối với tài khoản nền tảng.

    \item \textbf{Phí pilot ưu đãi hoặc khoản đặt cọc:}

    Audit ban đầu có thể miễn phí, nhưng pilot nên có phí hoặc khoản đặt cọc nhỏ. Khoản này có thể được khấu trừ vào gói tiếp theo nếu khách hàng tiếp tục sử dụng.

    Việc thu một khoản phí hoặc đặt cọc giúp nhóm phân biệt giữa sự quan tâm mang tính tham khảo và nhu cầu có mức cam kết thực tế.
\end{enumerate}

Nhóm không sử dụng hình thức ``miễn phí không giới hạn'' vì phương án này không kiểm chứng được mức sẵn sàng chi trả và có thể làm phát sinh khối lượng hỗ trợ vượt quá nguồn lực.

\subsubsection{Điều kiện tham gia pilot}

Một đơn vị chỉ được xác nhận tham gia pilot khi đáp ứng đầy đủ các điều kiện sau:

\begin{enumerate}
    \item thuộc nhóm khách hàng mục tiêu đã được xác định;
    \item có chiến dịch, khóa học hoặc nhu cầu livestream cụ thể;
    \item có bộ dữ liệu tối thiểu phục vụ cấu hình;
    \item có quyền sử dụng tài khoản, hình ảnh, video, giọng nói và nội dung cung cấp;
    \item chỉ định người phụ trách dự án và người duyệt thông tin;
    \item thống nhất phạm vi, thời gian và tiêu chí đánh giá pilot;
    \item chấp nhận rằng sản phẩm đang ở giai đoạn MVP và có thể phát sinh lỗi;
    \item đồng ý cho nhóm ghi nhận dữ liệu vận hành cần thiết;
    \item đồng ý tham gia buổi tổng kết và đưa ra quyết định sau pilot;
    \item hoàn tất khoản phí hoặc đặt cọc nếu pilot áp dụng hình thức có thu phí.
\end{enumerate}

Nếu khách hàng thiếu dữ liệu hoặc chưa có người phụ trách, nhóm có thể hỗ trợ chuẩn bị trước nhưng chưa bắt đầu pilot. Ngày bắt đầu chỉ được xác nhận sau khi các điều kiện tối thiểu đã hoàn tất.

\subsubsection{Quy trình từ tiếp cận đến pilot}

Quy trình thu hút khách hàng đầu tiên gồm tám bước.

\begin{table}[htbp]
    \centering
    \caption{Quy trình từ tiếp cận đến triển khai pilot}
    \label{tab:first-customer-to-pilot-process}
    \begin{tabular}{
        p{0.07\textwidth}
        p{0.20\textwidth}
        p{0.43\textwidth}
        p{0.19\textwidth}
    }
        \hline
        \textbf{Bước}
        & \textbf{Giai đoạn}
        & \textbf{Hoạt động}
        & \textbf{Đầu ra bắt buộc} \\
        \hline

        1
        &
        Tìm kiếm
        &
        Thu thập thông tin từ mạng lưới, Google Maps, Facebook, TikTok, cộng đồng và đối tác.
        &
        Hồ sơ khách hàng tiềm năng. \\

        \hline

        2
        &
        Sàng lọc
        &
        Chấm điểm theo phân khúc, nhu cầu, dữ liệu, quy trình lead và khả năng tiếp cận người quyết định.
        &
        Điểm ICP và mức ưu tiên. \\

        \hline

        3
        &
        Tiếp cận
        &
        Gửi tin nhắn cá nhân hóa, nêu vấn đề quan sát được và đề nghị audit hoặc trao đổi.
        &
        Phản hồi hoặc lịch trao đổi. \\

        \hline

        4
        &
        Khám phá nhu cầu
        &
        Phỏng vấn về quy trình hiện tại, tần suất livestream, dữ liệu, nhân sự, khó khăn và mục tiêu.
        &
        Biên bản nhu cầu và điều kiện thành công. \\

        \hline

        5
        &
        Audit và demo
        &
        Phân tích một tình huống thực tế và trình diễn workflow phù hợp với dữ liệu của khách hàng.
        &
        Báo cáo audit và phản hồi demo. \\

        \hline

        6
        &
        Đề xuất pilot
        &
        Thống nhất phạm vi, trách nhiệm, dữ liệu, thời gian, chi phí và chỉ số đánh giá.
        &
        Đề xuất pilot được chấp thuận. \\

        \hline

        7
        &
        Onboarding
        &
        Nhận dữ liệu, chuẩn hóa, thiết lập quy tắc, kiểm thử và duyệt trước khi vận hành.
        &
        Checklist onboarding hoàn tất. \\

        \hline

        8
        &
        Triển khai và đánh giá
        &
        Chạy phiên thử nghiệm, ghi nhận dữ liệu, xử lý lỗi và tổ chức buổi tổng kết.
        &
        Báo cáo pilot và quyết định tiếp theo. \\

        \hline
    \end{tabular}
\end{table}

\subsubsection{Tiêu chí hoàn tất pilot}

Pilot được xem là hoàn tất khi:

\begin{itemize}
    \item bộ dữ liệu đã được khách hàng phê duyệt;
    \item ít nhất một phiên thử nghiệm được triển khai thành công;
    \item hệ thống ghi nhận được log xử lý câu hỏi và các trường hợp chuyển người thật;
    \item thông tin khách hàng tiềm năng được chuyển sang quy trình tư vấn đã thống nhất;
    \item các lỗi nghiêm trọng đã được ghi nhận và xử lý;
    \item khách hàng tham gia buổi đánh giá cuối kỳ;
    \item khách hàng đưa ra quyết định tiếp tục, điều chỉnh hoặc dừng cùng với lý do cụ thể.
\end{itemize}

Một phiên chỉ được xem là thành công về mặt kỹ thuật khi chạy đúng phạm vi đã thống nhất. Điều này chưa đồng nghĩa với việc pilot đã chứng minh được hiệu quả kinh doanh.

\subsubsection{Điều kiện chuyển sang trả phí}

Sau khi pilot kết thúc, khách hàng được đề nghị chuyển sang gói trả phí khi đồng thời xuất hiện các điều kiện sau:

\begin{enumerate}
    \item \textbf{Sản phẩm tạo ra giá trị vận hành có thể quan sát:}

    Đội ngũ khách hàng xác nhận rằng một phần công việc lặp lại đã được giảm, quy trình rõ ràng hơn hoặc khả năng theo dõi phiên được cải thiện.

    \item \textbf{Chất lượng nằm trong ngưỡng chấp nhận:}

    Nội dung, phản hồi và quy trình chuyển người thật đáp ứng tiêu chuẩn mà khách hàng đã thống nhất trước pilot.

    \item \textbf{Onboarding có thể lặp lại:}

    Khách hàng có khả năng tiếp tục cung cấp và cập nhật dữ liệu mà không phụ thuộc hoàn toàn vào nhóm 1Stream.

    \item \textbf{Có nhu cầu sử dụng tiếp:}

    Trung tâm có lịch tuyển sinh, chiến dịch hoặc kế hoạch livestream trong tháng tiếp theo.

    \item \textbf{Người quyết định chấp thuận ngân sách:}

    Chủ trung tâm hoặc người có quyền phê duyệt xác nhận mức chi phí phù hợp với giá trị nhận được.

    \item \textbf{Hai bên thống nhất phạm vi dịch vụ:}

    Số kênh, số giờ, giới hạn nội dung, mức hỗ trợ và trách nhiệm của mỗi bên được xác định rõ.
\end{enumerate}

Các chỉ số hỗ trợ quyết định chuyển đổi gồm:

\begin{itemize}
    \item thời gian từ khi nhận đủ dữ liệu đến phiên đầu tiên;
    \item số lượng và tỷ lệ câu hỏi phổ biến được xử lý trong phạm vi cho phép;
    \item tỷ lệ trường hợp ngoài phạm vi được chuyển đúng cho nhân viên;
    \item số khách hàng tiềm năng được ghi nhận và bàn giao;
    \item thời gian phản hồi trung bình;
    \item số giờ công việc mà nhân viên tự đánh giá đã được giảm;
    \item mức hài lòng của người sử dụng trực tiếp;
    \item nhu cầu tiếp tục sử dụng trong chiến dịch tiếp theo.
\end{itemize}

Nhóm không sử dụng một mức tăng doanh thu cụ thể làm điều kiện bắt buộc trong giai đoạn đầu, vì doanh thu tuyển sinh còn chịu ảnh hưởng bởi nhiều yếu tố như chất lượng khóa học, thương hiệu, ưu đãi, ngân sách quảng cáo và năng lực tư vấn của trung tâm.

\subsubsection{Các trường hợp chưa chuyển sang trả phí}

Khách hàng chưa được đề nghị chuyển sang gói trả phí nếu:

\begin{itemize}
    \item chưa hoàn tất phiên thử nghiệm;
    \item dữ liệu đầu vào thiếu hoặc thường xuyên sai;
    \item nội dung chưa được người có thẩm quyền phê duyệt;
    \item có lỗi nghiêm trọng liên quan đến học phí, lịch học hoặc thông tin cam kết;
    \item chưa xác định được quy trình nhận và xử lý lead;
    \item khách hàng chưa có nhu cầu sử dụng trong thời gian tiếp theo;
    \item giá trị tạo ra chưa đủ rõ so với quy trình hiện tại;
    \item người quyết định chưa tham gia đánh giá kết quả.
\end{itemize}

Trong trường hợp này, nhóm cần xác định nguyên nhân đến từ sản phẩm, quy trình onboarding, thông điệp, khách hàng không phù hợp hay thời điểm triển khai chưa phù hợp. Kết quả được sử dụng để điều chỉnh chiến lược thay vì tiếp tục bán hàng bằng mọi giá.

\subsubsection{Chiến lược tạo khách hàng tiếp theo từ khách hàng đầu tiên}

Khi một pilot tạo ra kết quả tích cực, nhóm sử dụng kết quả đó để tạo vòng lặp thu hút khách hàng mới.

Quy trình vòng lặp được xác định như sau:

\begin{center}
\textbf{Khách hàng tham gia pilot}
$\rightarrow$
\textbf{Thu thập dữ liệu và phản hồi}
$\rightarrow$
\textbf{Xây dựng case study}
$\rightarrow$
\textbf{Nhận lời giới thiệu}
$\rightarrow$
\textbf{Tiếp cận đơn vị tương tự}
$\rightarrow$
\textbf{Triển khai pilot mới}
\end{center}

Các hình thức khai thác bằng chứng gồm:

\begin{itemize}
    \item case study có số liệu đã được khách hàng đồng ý;
    \item phản hồi hoặc trích dẫn của người sử dụng;
    \item video mô tả quy trình trước và sau pilot;
    \item báo cáo tổng hợp đã ẩn danh;
    \item lời giới thiệu đến một trung tâm khác;
    \item workshop nhỏ có sự tham gia của khách hàng pilot.
\end{itemize}

Nhóm chỉ công bố tên, logo, số liệu hoặc phát biểu của khách hàng khi có sự đồng ý. Nếu khách hàng không đồng ý công khai, nhóm chỉ sử dụng dữ liệu tổng hợp và ẩn danh.

\subsubsection{Chỉ số đánh giá chiến lược khách hàng đầu tiên}

\begin{table}[htbp]
    \centering
    \caption{Chỉ số đánh giá chiến lược thu hút khách hàng đầu tiên}
    \label{tab:first-customer-acquisition-kpis}
    \begin{tabular}{
        p{0.31\textwidth}
        p{0.34\textwidth}
        p{0.23\textwidth}
    }
        \hline
        \textbf{Chỉ số}
        & \textbf{Cách tính}
        & \textbf{Mục tiêu thử nghiệm} \\
        \hline

        Tỷ lệ phản hồi
        &
        Số phản hồi có ý nghĩa chia cho số đơn vị được tiếp cận.
        &
        Tối thiểu 25\% \\

        \hline

        Tỷ lệ chuyển sang demo
        &
        Số cuộc demo hoàn tất chia cho số khách hàng phản hồi.
        &
        Tối thiểu 50\% \\

        \hline

        Tỷ lệ gửi dữ liệu mẫu
        &
        Số khách hàng gửi dữ liệu chia cho số demo hoàn tất.
        &
        Tối thiểu 40\% \\

        \hline

        Tỷ lệ chuyển sang pilot
        &
        Số pilot bắt đầu chia cho số demo hoàn tất.
        &
        Tối thiểu 40\% \\

        \hline

        Tỷ lệ hoàn tất onboarding
        &
        Số đơn vị hoàn tất onboarding chia cho số pilot bắt đầu.
        &
        Tối thiểu 80\% \\

        \hline

        Tỷ lệ chuyển sang trả phí
        &
        Số đơn vị trả phí hoặc đặt cọc chia cho số pilot hoàn tất.
        &
        Tối thiểu 1 trong 3 đơn vị \\

        \hline

        Tỷ lệ giới thiệu
        &
        Số khách hàng giới thiệu đơn vị khác chia cho số pilot hoàn tất.
        &
        Tối thiểu 1 lượt giới thiệu \\

        \hline
    \end{tabular}
\end{table}

Do số lượng khách hàng trong giai đoạn đầu còn nhỏ, các tỷ lệ trên cần được đọc cùng với phản hồi định tính. Một khách hàng không tiếp tục nhưng cung cấp lý do rõ ràng vẫn tạo ra dữ liệu có giá trị cho việc điều chỉnh sản phẩm và chiến lược tiếp cận.

\subsubsection{Kết luận}

Chiến lược thu hút khách hàng đầu tiên của 1Stream được xây dựng theo hướng tiếp cận trực tiếp, có chọn lọc và dựa trên bằng chứng. Nhóm ưu tiên các trung tâm ngoại ngữ có nhu cầu livestream thật, dữ liệu tương đối chuẩn hóa, đội ngũ vận hành mỏng và có người phụ trách phối hợp.

Khách hàng được tìm từ mạng lưới giới thiệu, Google Maps, Facebook, TikTok, cộng đồng giáo dục, đối tác dịch vụ và các sự kiện chuyên ngành. Audit miễn phí được sử dụng để tạo giá trị ban đầu, trong khi pilot được giới hạn về khóa học, nền tảng, thời gian và phạm vi chức năng nhằm giảm rủi ro cho cả hai bên.

Việc chuyển sang trả phí không dựa trên lời khen hoặc số lượt xem, mà dựa trên khả năng hoàn tất onboarding, chất lượng vận hành, nhu cầu sử dụng tiếp, mức độ chấp nhận của đội ngũ và quyết định ngân sách của người có thẩm quyền. Cách tiếp cận này giúp 1Stream kiểm chứng nhu cầu thị trường bằng khách hàng thật trước khi mở rộng quảng bá hoặc đầu tư ngân sách lớn hơn.